\ExerciseSection

\begin{enumerate}
\def\labelenumi{\arabic{enumi})}
\item
  Give a proof of the fact that \(\mathbf{R}^n\) is equipotent with \(\mathbf{R}\) by using Cantor's theorem (Chapter IV, § 8, no. 6, Theorem 1) and the relation \(2^{\mathfrak{a}} \cdot 2^{\mathfrak{a}} = 2^{\mathfrak{a}}\) , valid for all infinite cardinals \(\mathfrak{a}\) .
\item
  There exists a continuous mapping of the interval \(\mathbf{I} = [\mathrm{o},\mathrm{i}]\) of \(\mathbf{R}\) onto the square \(\mathbf{I}\times \mathbf{I}\) of \(\mathbf{R}^2\) (the ``Peano curve'', cf.~Exercise 8). (Show first, with the help of Exercise 11 of Chapter IV, § 8, that there exists a continuous mapping \(f\) of Cantor's triadic set \(\mathbf{K}\) onto \(\mathbf{I}\times \mathbf{I}\) , and then extend \(f\) to \(\mathbf{I}\) ).
\item
  Let A and B be two countable dense subsets of \(R^{2}\) . Show that there exists a homeomorphism of \(R^{2}\) onto itself which maps A onto B. (Show first that, by means of a rotation, we can reduce to the case where the projections \(pr_{1}\) and \(pr_{2}\) are injective mappings of A and B into R. Then define, by a suitable inductive process, a bijection of A onto B which determines a monotone mapping of \(pr_{1}A\) onto \(pr_{1}B\) and a monotone mapping of \(pr_{2}A\) onto \(pr_{2}B\) . Finally, using Exercise 11, Chapter IV, § 2, show that this mapping is a homeomorphism which can be extended to a homeomorphism of \(R^{2}\) onto itself.) Deduce that, if C is a subset of \(R^{2}\) whose complement is dense, C is homeomorphic to a subset of the complement of \(Q^{2}\) in \(R^{2}\) . Generalize these results to \(R^{n}\) , n \textgreater{} 2.
\item
  Every countable subspace of \(\mathbf{R}^n\) with no isolated points is homeomorphic to the rational line \(\mathbf{Q}\) (apply Exercise 13 of Chapter IV, § 8).
\item
  Every totally discontinuous compact subspace of \(\mathbf{R}^n\) with no isolated points is homeomorphic to Cantor's triadic set (use Chapter IV, § 8, Exercises 11 and 12, and Chapter II, § 4, no. 4, Proposition 6).
\item
  A subset L of \(R^{n}\) is a broken line if there exists a finite sequence \((x_{i})_{0\leqslant i\leqslant p}\) of points of \(R^{n}\) such that, if \(S_{i}\) denotes the segment with end-points \(x_{i-1}\) and \(x_{i}\) ( \(i\leqslant i\leqslant p\) ), L is the union of the \(S_{i}\) , which are called the sides of L. (In general there is an infinite number of finite sequences of points of \(R^{n}\) which define the same broken line.) A broken line is also the image of a mapping u of {[}o, i{]} into \(R^{n}\) such that there exists a strictly increasing sequence \((t_{j})_{0\leqslant j\leqslant q}\) in {[}o, i{]} with \(t_{0}=0\) and \(t_{q}=1\) , with the property that u is an affine linear mapping
\end{enumerate}

\[
t \rightarrow \boldsymbol {a} _ {j} + t \boldsymbol {b} _ {j} \quad \text { for } \quad t _ {j - 1} \leqslant t \leqslant t _ {j} \quad \text { and } \quad i \leqslant j \leqslant q
\]

(such a mapping \(u\) is said to be ``piecewise linear''). Given a non-empty subset A of \(\mathbf{R}^n\) we say that two points \(a, b\) of A can be joined by a broken line in A if there exists a broken line \(L \subset A\) defined by a sequence \((x_i)_{0 \leqslant i \leqslant p}\) such that \(x_0 = a\) and \(x_p = b\) . If any two points of A can be joined by a broken line in A, then A is connected. Conversely, if A is a connected open subset of \(\mathbf{R}^n\) , show that any two points of A can be joined by a broken line in A (consider the relation `` \(x\) and \(y\) can be joined by a broken line in A'' between points \(x, y\) of A; show that this is an equivalence relation whose classes are open sets); we can even assume always that this broken line is a union of segments each of which is parallel to a coordinate axis (same method). Deduce that, if A is a non-empty open set in \(\mathbf{R}^n\) , the component of a point \(a \in A\) is the set of all points of A which can be joined to \(a\) by a broken line in A.

\begin{enumerate}
\def\labelenumi{\arabic{enumi})}
\setcounter{enumi}{6}
\item
  Let A be a non-empty connected open set in \(R^{n}\) (n \textgreater{} 1) and let \((\mathrm{V}_{p})_{p \in \mathbb{N}}\) be a countable family of linear varieties in \(R^{n}\) , each of which is of dimension \(\leqslant n - 2\) . Show that if B is the union of the \(V_{p}\) , then \(A \cap \mathfrak{C}B\) is dense in A and connected. (To show that \(A \cap \mathfrak{C}B\) is dense in A, use induction on n.~Then show that if x, y are two distinct points of A, there exists for each \(\varepsilon > 0\) a point \(y' \in A\) such that \(||y' - y|| \leqslant \varepsilon\) and such that the closed segment with end-points x and \(y'\) meets B only at x. Finally use Exercise 6 to show that \(A \cap \mathfrak{C}B\) is connected.)
\item
  Deduce from Exercise 7 that, if \(n > \mathbf{I}\) , a non-empty open subset of \(\mathbf{R}^n\) is not homeomorphic to a subset of \(\mathbf{R}\) (*).
\item
  A broken line L in \(\mathbf{R}^{n}\ (n > \mathrm{i})\) (Exercise 6) is said to be simple if it is homeomorphic to the interval \(I = [0, i]\) of R. It amounts to the same thing to say that there exists a piecewise linear homeomorphism u of I onto L (Exercise 6). Show that if A is a connected open set in \(\mathbf{R}^2\) and if \(\mathbf{L}\subset\mathbf{A}\) is a simple broken line, then \(\mathbf{A}-\mathbf{L}\) is connected (argue by induction on the number of segments which make up L, and use Exercise 6 of Chapter I, § 11).
\end{enumerate}

* 10) a) Let M be a finite set and let R \(x, y\) be a symmetric relation between elements x, y of M. Suppose that there exist two distinct elements a, b in M with the following properties: there exist x ≠ a and y ≠ b such that x (resp. y) is the only element z ∈ M such that R \(a, z\) (resp. R \(b, z\) ) is true; and for each t ∈ M other than a and b, the set of all z ∈ M such that R \(t, z\) is true is a set consisting of two elements, distinct from t. Show that under these conditions there exists a bijection i → x\_i of an interval {[}o, n{]} of N onto M such that x\_0 = a, x\_n = b and such that R \(x_{i-1}, x_i\) is true for i ≤ i ≤ n (define x\_i by induction on i).

\begin{enumerate}
\def\labelenumi{\alph{enumi})}
\setcounter{enumi}{1}
\tightlist
\item
  In \(\mathbf{R}^n\) ( \(n \geqslant 2\) ), identified with \(\mathbf{R}^{n-1} \times \mathbf{R}\) , let \(\mathbf{L}\) and \(\mathbf{L}'\) be two broken lines, defined respectively by two sequences \((\boldsymbol{x}_i)_{0 \leqslant i \leqslant p}\) and \((\boldsymbol{x}_j')_{0 \leqslant j \leqslant q}\) , with the following properties: (i) if we put \(\boldsymbol{x}_i = (\boldsymbol{y}_i, z_i)\) , \(\boldsymbol{x}_j' = (\boldsymbol{y}_j', z_j')\) , we have \(z_0 = z_0' = 0\) , \(z_p = z_q' = 1\) , \(0 < z_i < 1\) for \(1 \leqslant i \leqslant p - 1\) and \(0 < z_j' < 1\) for \(1 \leqslant j \leqslant q - 1\) ; (ii) the \(p + q - 2\) numbers \(z_i, z_j'\) ( \(1 \leqslant i \leqslant p - 1\) , \(1 \leqslant j \leqslant q - 1\) ) are all distinct; (iii) two distinct sides of \(\mathbf{L}\) (resp. \(\mathbf{L}'\) ) have at most one point in common {[}which may or may not be one of the \(\boldsymbol{x}_i\) (resp. \(\boldsymbol{x}_j'\) ){]}. Show that there exist two surjective piecewise linear mappings (Exercise 6) \(\boldsymbol{u}\) : \(\mathbf{I} \to \mathbf{L}\) , \(\boldsymbol{u}'\colon \mathbf{I} \to \mathbf{L}'\) , where \(\mathbf{I}\) is the interval \([0, 1]\) of \(\mathbf{R}\) , such that, if we put \(\boldsymbol{u}(t) = (\boldsymbol{v}(t), \zeta(t))\) , \(\boldsymbol{u}'(t) = (\boldsymbol{v}'(t), \zeta'(t))\) ( \(t \in \mathbf{I}\) ), we have \(\zeta(t) = \zeta'(t)\) for all \(t \in \mathbf{I}\) ,
\end{enumerate}

\[
\zeta (\mathrm{o}) = \zeta^ {\prime} (\mathrm{o}) = \mathrm{o} \quad \text { and } \quad \zeta (\mathrm{i}) = \zeta^ {\prime} (\mathrm{i}) = \mathrm{i}.
\]

{[}Let \((a_{k})_{1\leqslant k\leqslant p + q - 2}\) be the strictly increasing sequence formed by the \(z_{i}\) and the \(z_j^{\prime}\) other than o and i, and put \(a_0 = 0\) , \(a_{p + q - 1} = \mathrm{i}\) ; let \(\mathbf{B}_i\) denote the set of all \(\pmb {x} = (\pmb {y},z)\in \mathbb{R}^n\) such that

\[
a _ {i - 1} \leqslant z \leqslant a _ {i} \quad (\mathrm{I} \leqslant i \leqslant p + q - \mathrm{I});
\]

let \(\mathfrak{M}\) denote the set of all pairs \(\gamma = (\mathrm{C},\mathrm{C}^{\prime})\) , where \(\mathrm{C}\) (resp. \(\mathrm{C}'\) ) is the intersection of the same \(\mathbf{B}_i\) with a side of \(\mathbf{L}\) (resp. \(\mathbf{L}'\) ), such that neither \(\mathbf{C}\) nor \(\mathbf{C}'\) consists of a single point; and let \(\mathbf{R}\wr \alpha, \beta\wr\) denote the following relation between elements \(\alpha, \beta\) of \(\mathfrak{M}\) : `` \(\alpha \neq \beta\) ; if \(\alpha = (\mathrm{C}_1,\mathrm{C}_1')\) , \(\beta = (\mathrm{C}_2,\mathrm{C}_2')\) , then \(\mathrm{C}_1 \cap \mathrm{C}_2\) and \(\mathrm{C}_1' \cap \mathrm{C}_2'\) are not empty, one of them consists of a single point, and if this point is not one of the \(x_i\) (resp. \(x_j'\) ) then \(\mathrm{C}_1\) and \(\mathrm{C}_2\) (resp. \(\mathrm{C}_1'\) and \(\mathrm{C}_2'\) ) are both contained in the same side of \(\mathbf{L}\) (resp. \(\mathbf{L}'\) ). Show that \(a\) ) is applicable to the relation \(\mathbf{R}.\) {]}

\begin{enumerate}
\def\labelenumi{\alph{enumi})}
\setcounter{enumi}{2}
\tightlist
\item
  In \(\mathbf{R}^n\) , identified with \(\mathbf{R}^{n-1} \times \mathbf{R}\) , let \(\mathbf{K}\) be a connected compact set such that \(\mathbf{K} \cap (\mathbf{R}^{n-1} \times \{\mathbf{o}\})\) contains at least two distinct points. Show that if \(\mathbf{K}' = \mathbf{K} - \mathbf{K}\) (the set of all \(x - y\) , where \(x \in \mathbf{K}\) and \(y \in \mathbf{K}\) ), then \(\mathbf{K}' \cap (\mathbf{R}^{n-1} \times \{\mathbf{o}\})\) contains a connected set which does not consist of a single point. {[}Use b), applying Chapter II, § 4, Proposition 6 and Exercise 15.{]}
\end{enumerate}

\begin{enumerate}
\def\labelenumi{\arabic{enumi})}
\setcounter{enumi}{10}
\tightlist
\item
  Extend Exercises 14 and 15 of Chapter IV, § 8 to the spaces \(\mathbf{R}^n\) ( \(n \geqslant 2\) ).
\end{enumerate}

* 12) a) Let \(f\) be a mapping of \(\mathbf{R}\) into \(\mathbf{R}\) , and let \(\mathbf{G} \subset \mathbf{R}^2\) be the graph of \(f\) . Suppose that \(\mathbf{G}\) is dense in \(\mathbf{R}^2\) and meets every perfect set in \(\mathbf{R}^2\) which is not contained in any countable union of lines \(\{x_n\} \times \mathbf{R}\) . Show that \(\mathbf{G}\) is connected. (If this were false, show that there would exist two non-empty disjoint open sets \(\mathbf{A}\) , \(\mathbf{B}\) in \(\mathbf{R}^2\) such that \(\mathbf{G}\) was the union of \(\mathbf{G} \cap \mathbf{A}\) and \(\mathbf{G} \cap \mathbf{B}\) , and there would exist at least one point of \(\mathbf{A}\) and one point of \(\mathbf{B}\) with the same abscissa, by using the fact that \(\mathbf{R}\) is connected; then use Exercise 11.)

\begin{enumerate}
\def\labelenumi{\alph{enumi})}
\setcounter{enumi}{1}
\tightlist
\item
  Deduce from a) that there is a linear mapping f of R into R whose graph G is dense in \(R^{2}\) and connected. {[}Using a) and Exercise 11, define by transfinite induction the values of f at the points of a Hamel base of R, using the same method as in Set Theory, Chapter III, § 6, Exercise 24.{]} Deduce that the subgroup G of \(R^{2}\) satisfies conditions ( \(R_{I}\) ) and ( \(R_{II}\) ) of Chapter V, § 3, Exercise 4, but is not locally compact, and therefore has a topology strictly finer than the usual topology of R; and that G is not locally connected.
\end{enumerate}

\begin{enumerate}
\def\labelenumi{\arabic{enumi})}
\setcounter{enumi}{12}
\tightlist
\item
  Identifying \(R^{n}\) with \(R^{n-1} \times R\) , let f be a continuous mapping of an open subset A of \(R^{n-1}\) into R, and let S be its graph. Show that the subspace S of \(R^{n}\) is homeomorphic to A, and that the complement of S in the ``cylinder'' \(A \times R\) is a dense open set in \(A \times R\) , and has exactly two components if A is connected.
\end{enumerate}

